\section{Summary of Contributions}
This thesis developed and evaluated EQUITAS-RCMF, a multimodal architecture for stress detection designed to prevent a known spurious visual attribute, a facial scar, from influencing its predictions. Its main results are as follows.
\begin{enumerate}
    \item \textbf{Physiological information in facial video.} A pre-registered test found recoverable physiological information in the UBFC-Phys facial video (POS, $p = 0.0314$), with consistent exploratory results (CHROM, $p = 0.0001$; PBV, $p = 0.0103$). The visual branch therefore has access to both genuine and spurious cues.
    \item \textbf{Orthogonal subspace decomposition.} Visual features were projected onto orthogonal causal and confounder subspaces by a thin QR projection applied at every forward pass, with an orthogonality error of $1.7149 \times 10^{-6}$ at initialization.
    \item \textbf{Counterfactual invariance.} In the demographic audit of the seed-42 model, removing the scar changed predictions very little, with a counterfactual gap of 0.0004--0.0007 across subgroups.
    \item \textbf{Penalty alone is insufficient.} A model trained with a counterfactual penalty but without the subspace decomposition reduced the counterfactual gap to 0.0019 but reached only 50.20\% accuracy, close to chance, which motivates constraining the structure of the representation rather than relying on the loss alone.
    \item \textbf{Deployment.} The compiled model does not require the scar mask and runs at 0.90~ms per frame (1112.7 frames per second) on an x86-64 CPU.
\end{enumerate}

\section{Boundaries of the Results}
\subsection{Cohort Size}
The eight-participant cohort is a demanding setting for shortcut learning, because a network with over a million parameters, trained on 2,344 windows from a handful of people, can easily memorize a spurious feature. It also limits what can be concluded: subgroup estimates rest on very few participants, and significance testing was possible only at the clip level. Evaluation on the full UBFC-Phys cohort of 56 participants and on data from other sites is needed before the results can be generalized.

\subsection{Accuracy and Fairness}
Subgroup accuracies ranged from about 59.7\% to 64.3\%, while the counterfactual gap stayed below 0.001. The low counterfactual gap is the main result, but it must be interpreted together with accuracy: the model shows modality collapse, with accuracies below the approximately 72.15\% of the physiology-only model, so adding the visual modality currently lowers accuracy rather than improving it (Section~\ref{sec:limitations}). Achieving counterfactual invariance while retaining the benefit of the visual modality is the central open problem for this architecture.
% TODO(authors): If the code check in Chapter 3 confirms that the privileged Focus
% (log(1 + r), unbounded) and Focus_auto (sigmoid, in (0, 1)) are on different scales,
% add a subsection here explaining that the deployed gate may penalize scar-focused
% visual features less strongly than during training.

\subsection{Quantization and Orthogonality}
Standard post-training quantization rounds weights to a low-precision grid, which would perturb the projection matrix so that its rows are no longer exactly orthonormal. The effect of this perturbation on the separation between the causal and confounder subspaces has not been measured. Quantization methods that preserve orthogonality, or quantization-aware training with an orthogonality constraint, are therefore a necessary direction before the model is deployed in reduced precision.

\section{Future Work}
\subsection{Evaluation}
The most immediate extensions are to recover the master checkpoint and report the pending sanity-check results; to report accuracy and fairness for every test regime, mode, and seed; to evaluate with participant-level splits; and to compare EQUITAS-RCMF with a vision-only variant and with group-robust baselines that use the known scar groups~\cite{sagawa2020}.
% TODO(authors): Delete "to evaluate with participant-level splits" if the split
% used in this thesis is already by participant (see the Chapter 3 TODO).

\subsection{Addressing Modality Collapse}
The first step toward resolving modality collapse is diagnostic: measuring the gate values for scarred and unscarred inputs, testing accuracy when the visual input is removed, and checking whether the focus signal has the same scale during training, when it is computed from the scar mask, as at inference, when it is estimated from the confounder subspace. A mismatch between the two would cause the deployed gate to behave differently from the trained one. Using the original, pre-scar frame as the counterfactual, where the code does not already do so, would also make the counterfactual penalty exact, since the two inputs would then differ only in the scar.

\subsection{Alternative Backbones}
The subspace decomposition does not depend on a particular backbone. Vision Mamba replaces self-attention with a bidirectional state-space model whose cost grows linearly rather than quadratically with the number of image patches~\cite{zhu2024vim}, which could make higher-resolution facial crops feasible on constrained hardware, although efficient implementations of its scan operation currently target GPUs, and its advantage on mobile CPUs has not been established. RepViT incorporates design elements of vision transformers into a mobile convolutional network through structural re-parameterization, with a reported ImageNet top-1 accuracy above 80\% at about 1.0~ms latency on an iPhone~12~\cite{wang2024repvit}. Replacing the MobileNetV3-Small backbone with either architecture would test whether the approach generalizes across backbone designs.

\section{Concluding Remarks}
This thesis set out to design an edge-deployable model whose stress predictions do not depend on a known spurious visual attribute. The results support part of that aim. EQUITAS-RCMF runs well within real-time limits, does not require the scar annotation at deployment, and in the reported audit its predictions change very little when the scar is removed. The results also show why these findings must be read with care: the facial video carries genuine physiological information, yet the fused model is currently less accurate than a model that uses physiology alone. Counterfactual invariance to the scar has therefore so far been achieved alongside, rather than in addition to, a loss of the benefit the visual modality should provide. Separating these two effects, through the pending per-regime results, gate analysis, and vision-ablation tests, and extending the evaluation to larger cohorts, will determine whether structural constraints of this kind can make multimodal biometric systems both fair and useful.
